\section*{Capítulo \ref{ch:b2:chemical-potential} --- Potencial químico y mezclas}

\begin{solution}{exo:b2:chemical-potential:1}
La tangente en $x_2 = 0.4$ corta los bordes en $\bar V_1 \approx 14$ y $\bar
V_2 \approx 49$ (unidades del modelo). Entonces $0.6 \times 14 + 0.4 \times 49 = 28.0$,
el volumen molar leído en la curva en $x_2 = 0.4$, por debajo de los
$0.6 \times 18 + 0.4 \times 58 = 34.0$ de los volúmenes puros.
\end{solution}

\begin{solution}{exo:b2:chemical-potential:2}
$n(\text{agua}) = 1000/18.015 = \qty{55.5}{mol}$, $x_1 = 55.5/56.5 = 0.982$;
$p = 0.982 \times 3.17 = \qty{3.11}{kPa}$: un descenso del 1.8\,\%.
\end{solution}

\begin{solution}{exo:b2:chemical-potential:3}
$p(\ce{O2}) = 0.2095 \times \qty{1.013e5}{Pa} = \qty{2.12e4}{Pa}$;
$c = 1.3 \times 10^{-5} \times 2.12 \times 10^4 = \qty{0.28}{mol/m^3} =
\qty{2.8e-4}{mol/L}$, es decir, $2.8 \times 10^{-4} \times 32.0 \times 10^3 =
\qty{8.8}{mg/L}$.
\end{solution}

\begin{solution}{exo:b2:chemical-potential:4}
$Q = (1.0)^2/[(1.0)(0.010)^3] = 1.0 \times 10^6$; $RT\ln Q = 2.479 \times 13.82 =
\qty{34.2}{kJ/mol}$; $\Delta_r G = -32.8 + 34.2 = \qty{+1.4}{kJ/mol} > 0$: en esta
mezcla, tan pobre en hidrógeno, el amoníaco se descompone.
\end{solution}

\begin{solution}{exo:b2:chemical-potential:5}
Gibbs--Duhem: $x_1\dd\bar V_1 + x_2\dd\bar V_2 = 0$, luego $\dd\bar V_2 =
-(0.70/0.30) \times (-0.30) = \qty{+0.70}{mL/mol}$.
\end{solution}

\begin{solution}{exo:b2:chemical-potential:6}
$c(\text{partículas}) = 2 \times 9.0/58.44 = \qty{0.308}{mol/L} =
\qty{308}{mol/m^3}$; $\Pi = 308 \times 8.314 \times 310 = \qty{7.9e5}{Pa}$,
unos \qty{7.9}{bar}. El interior de un glóbulo rojo tiene la misma \omterm{def:b2:chemical-potential:osmotic-pressure}{presión osmótica}:
no hay flujo neto de agua. En agua pura, el agua entra en la célula, que se hincha y
revienta.
\end{solution}

\begin{solution}{exo:b2:chemical-potential:7}
$c = \Pi/RT = 825/(8.314 \times 298.15) = \qty{0.333}{mol/m^3}$; en
\qty{1.000e-4}{m^3}, $n = \qty{3.33e-5}{mol}$, $M = 2.00/3.33 \times 10^{-5}
= \qty{6.0e4}{g/mol}$ (\qty{60}{kg/mol}). $\Delta T_f = 1.86 \times 3.33
\times 10^{-5}/0.100 = \qty{6.2e-4}{K}$: inmedible, mientras que la \omterm{def:b2:chemical-potential:osmotic-pressure}{presión
osmótica} es una columna de \qty{8}{cm} de agua.
\end{solution}

\begin{solution}{exo:b2:chemical-potential:8}
(a) $I = \qty{1.0e-3}{mol/kg}$, $\log\gamma_\pm = -0.51 \times 0.0316$,
$\gamma_\pm = 0.964$. (b) $I = \frac12(4 \times 1.0 + 1 \times 2.0) \times
10^{-3} = \qty{3.0e-3}{mol/kg}$, $\log\gamma_\pm = -0.51 \times 2 \times 0.0548$,
$\gamma_\pm = 0.879$. (c) $I = \frac12(4 + 4) \times 10^{-3} = \qty{4.0e-3}{mol/kg}$,
$\log\gamma_\pm = -0.51 \times 4 \times 0.0632$, $\gamma_\pm = 0.74$.
\end{solution}

\begin{solution}{exo:b2:chemical-potential:9}
$b = 0.310/1.86 = \qty{0.167}{mol/kg}$; $n = 0.167 \times 0.0500 =
\qty{8.33e-3}{mol}$; $M = 1.50/8.33 \times 10^{-3} = \qty{180}{g/mol}$ (una
hexosa, por ejemplo).
\end{solution}

\begin{solution}{exo:b2:chemical-potential:10}
$x_1\,\dd(Ax_2^2) + x_2\,\dd(Ax_1^2) = 2Ax_1x_2(\dd x_2 + \dd x_1) = 0$ ya que
$x_1 + x_2 = 1$. Cuando $x_1 \to 0$, $\gamma_1 \to \mathrm e^{A}$ y $p_1 =
\gamma_1x_1p_1^* \approx \mathrm e^{A}p_1^*\,x_1$: $k_{H,1} = p_1^*\mathrm e^A$.
Para $A = 1.1$, $\mathrm e^{1.1} = 3.0$: la pendiente de Henry triplica la de Raoult.
\end{solution}

\begin{solution}{exo:b2:chemical-potential:11}
(a) Un soluto volátil aporta su propia presión de vapor; el vapor ya no es
disolvente puro y la deducción (vapor puro del disolvente en
equilibrio con la disolución) deja de valer. (b) Si el soluto entra en el sólido, el
\omterm{def:b2:chemical-potential:chemical-potential}{potencial químico} del sólido también desciende, en $RT\ln x_1^{\text{s}}$. Si la
disolución sólida tiene la misma composición que el líquido, ambas rectas descienden
lo mismo y el punto de congelación no cambia: el descenso
exige un soluto que el cristal rechace.
\end{solution}

\begin{solution}{exo:b2:chemical-potential:12}
Crioscopia: $b_{\min} = 0.01/1.86 = \qty{5.4e-3}{mol/kg}$; ebulloscopia:
$0.01/0.513 = \qty{1.9e-2}{mol/kg}$. A $\qty{5.4e-3}{mol/L} =
\qty{5.4}{mol/m^3}$, $\Pi = 5.4 \times 8.314 \times 298 = \qty{1.3e4}{Pa}$,
más de un metro de agua; un osmómetro aprecia unos pocos pascales. La osmometría es, con mucho,
la más sensible, y por eso mide las grandes masas molares de los
polímeros, cuyas molalidades son minúsculas.
\end{solution}

\begin{solution}{pb:b2:chemical-potential:1}
\textbf{1.} $\mu_1 = \mu_1^*(T, p) + RT\ln x_1$.
\textbf{2.} Diol $30/62.07 = \qty{0.483}{mol}$, agua $70/18.015 =
\qty{3.886}{mol}$; $x_1 = 3.886/4.369 = 0.889$.
\textbf{3.} $p = 0.889 \times 3.17 = \qty{2.82}{kPa}$.
\textbf{4.} $b = 0.483/0.070 = \qty{6.90}{mol/kg}$.
\textbf{5.} El diol hierve a \qty{197}{\celsius}: su presión de vapor a temperatura
ambiente es muy inferior a la del agua.
\textbf{6.} $\mu_1^{\text{s}}(T_f) = \mu_1^{*\text{l}}(T_f) + RT_f\ln x_1$.
\textbf{7.} $\ln x_1 = -\Delta_{\text{fus}}G(T)/(RT)$ con
$\Delta_{\text{fus}}G = \mu_1^{*\text{l}} - \mu_1^{\text{s}}$, nula en $T_f^*$;
Gibbs--Helmholtz, $\dd(\Delta_{\text{fus}}G/T)/\dd T = -\Delta_{\text{fus}}H/T^2$,
integrada de $T_f^*$ a $T_f$, da $\Delta_{\text{fus}}G(T_f)/T_f =
\Delta_{\text{fus}}H(1/T_f - 1/T_f^*)$, de donde se sigue la relación.
\textbf{8.} $\ln x_1 \approx -x_2 \approx -bM_1$ y $1/T_f - 1/T_f^* \approx
-\Delta T_f/T_f^{*2}$: $\Delta T_f = (RT_f^{*2}M_1/\Delta_{\text{fus}}H)\,b$.
\textbf{9.} $333.4 \times 18.015 = \qty{6007}{J/mol}$; $K_f = 8.314 \times
273.15^2 \times 0.018015/6007 = \qty{1.86}{K.kg/mol}$.
\textbf{10.} Que $\Delta_{\text{fus}}H$ no varía entre $T_f$ y
$T_f^*$ (sí varía, por Kirchhoff, alrededor de un 2\,\% cada \qty{3}{K} aquí).
\textbf{11.} $\Delta T_f = 1.86 \times 6.90 = \qty{12.8}{K}$: \qty{-12.8}{\celsius}.
\textbf{12.} $1/T_f = 1/273.15 - (8.314/6007)\ln 0.889$: $T_f = \qty{261.6}{K}$,
\qty{-11.6}{\celsius}.
\textbf{13.} La relación ideal exacta, que no supone una disolución
diluida (aquí el 11\,\% de las moléculas son de diol). Ambas suponen una \omterm{def:b2:chemical-potential:ideal-mixture}{mezcla ideal}
y una entalpía de fusión constante; el agua y el diol, unidos por enlaces de hidrógeno,
no forman una \omterm{def:b2:chemical-potential:ideal-mixture}{mezcla ideal}, de modo que el punto de congelación real es otro.
\textbf{14.} Hielo puro: el diol se queda en el líquido.
\textbf{15.} Retirar agua en forma de hielo rebaja $x_1$ en el líquido y, según la
relación, la temperatura de equilibrio desciende: el refrigerante se congela en un intervalo
de temperaturas, no en un punto.
\textbf{16.} $K_b = 8.314 \times 373.12^2 \times 0.018015/40\,650 =
\qty{0.513}{K.kg/mol}$.
\textbf{17.} $\Delta T_b = 0.513 \times 6.90 = \qty{3.5}{K}$.
\textbf{18.} $1/T = 1/373.12 - (8.314/40\,650)\ln(2.0/1.013)$: $T = \qty{393.5}{K}$,
unos \qty{120}{\celsius}.
\textbf{19.} El tapón: \qty{20}{K} frente a \qty{3.5}{K} del diol.
\textbf{20.} $x_1 = \exp[-(6007/8.314)(1/253.15 - 1/273.15)] = 0.811$.
\textbf{21.} $x_2 = 0.189$: $w = 0.189 \times 62.07/(0.189 \times 62.07 + 0.811
\times 18.015) = 0.44$.
\textbf{22.} Ley de las disoluciones diluidas: $b = 20/1.86 = \qty{10.75}{mol/kg}$, $w = 10.75 \times
62.07/(1000 + 10.75 \times 62.07) = 0.40$. Las linealizaciones ($\ln x_1 \approx
-x_2$, $x_2 \approx bM_1$) exageran el descenso por mol de diol a una concentración
tan alta, así que la ley de las disoluciones diluidas pide menos diol.
\textbf{23.} En el modelo de \omterm{def:b2:chemical-potential:ideal-mixture}{mezcla ideal}, un refrigerante con una fracción másica de
diol $\boldsymbol{w \approx 0.44}$ sigue líquido hasta \qty{-20}{\celsius}.
\end{solution}
